% ============================================
% CHAPTER 2 - PART 2: Phát hiện bất thường
% ============================================
% Tiếp tục từ section 2.3 (đã có 2.1-2.2)
% Task 2.5: Viết các section về anomaly detection methods

\section{Phát hiện bất thường trên dữ liệu log}
\label{sec:anomaly-detection}

\subsection{Định nghĩa và phân loại bất thường}
\label{subsec:anomaly-definition}

Trong bối cảnh phát hiện bất thường trên log hệ thống, cần phân biệt ba khái niệm thường bị nhầm lẫn: \textit{Anomaly} (bất thường), \textit{Fault} (lỗi), và \textit{Failure} (hư hỏng) \cite{chandola2009anomaly}.

\begin{itemize}
    \item \textbf{Anomaly (Bất thường):} Sự chênh lệch so với hành vi mong đợi, quan điểm từ dữ liệu.
    \item \textbf{Fault (Lỗi):} Nguyên nhân gốc rễ của vấn đề, quan điểm từ hệ thống.
    \item \textbf{Failure (Hư hỏng):} Sự cố có thể quan sát được, quan điểm từ người dùng.
\end{itemize}

\textbf{Ví dụ minh họa:} Một \textit{fault} (memory leak trong service X) có thể tạo ra \textit{anomaly} (tần suất log "OutOfMemory" tăng cao), cuối cùng dẫn đến \textit{failure} (service crash, người dùng thấy lỗi 503).

Phát hiện bất thường trên log (LogAD) tập trung vào việc phát hiện \textit{anomaly} như một \textit{proxy} để dự đoán \textit{fault} trước khi chúng gây ra \textit{failure}.

\subsubsection{Các loại bất thường}

Theo Chandola et al. \cite{chandola2009anomaly}, bất thường được phân loại thành ba loại:

\paragraph{Point Anomaly:} Một sự kiện đơn lẻ bất thường.

\textbf{Ví dụ:} Một dòng log chưa từng xuất hiện trong quá trình training:
\begin{verbatim}
Normal:   INFO: Request processed in 20ms
Normal:   INFO: Request processed in 25ms
Anomaly:  FATAL: Segmentation fault at address 0x0000
Normal:   INFO: Request processed in 22ms
\end{verbatim}

\paragraph{Collective Anomaly (Sequential Anomaly):} Một chuỗi sự kiện bất thường khi xét tổng thể, dù các sự kiện đơn lẻ có thể bình thường.

\textbf{Ví dụ:} Chuỗi bình thường là A→B→C, nhưng chuỗi A→C→B là bất thường (bỏ qua authentication).

\begin{table}[ht]
\centering
\caption{So sánh các loại bất thường}
\label{tab:anomaly-types}
\begin{tabular}{lccl}
\hline
\textbf{Loại} & \textbf{Sự kiện đơn} & \textbf{Chuỗi} & \textbf{Phương pháp phát hiện} \\
\hline
Point & Bất thường & N/A & Frequency-based, Outlier detection \\
Collective & Bình thường & Bất thường & Sequence modeling (LSTM, Transformer) \\
Contextual & Phụ thuộc & Phụ thuộc & Conditional models \\
\hline
\end{tabular}
\end{table}

\paragraph{Contextual Anomaly:} Một sự kiện chỉ bất thường trong ngữ cảnh cụ thể.

\textbf{Ví dụ:} 10,000 login requests là bình thường vào giờ hành chính (9h-18h), nhưng bất thường vào lúc nửa đêm (có thể là tấn công).

\textbf{Phạm vi nghiên cứu:} Luận văn này tập trung vào \textit{Collective Anomaly} vì đây là loại phổ biến nhất trong log hệ thống \cite{he2020loghub}.

\subsection{Các hướng tiếp cận theo mức độ giám sát}
\label{subsec:supervision-approaches}

\subsubsection{Supervised Learning (Học có giám sát)}

Huấn luyện trên dữ liệu đã gán nhãn (cả normal và anomaly).

\textbf{Ưu điểm:}
\begin{itemize}
    \item Độ chính xác cao nếu nhãn chính xác
    \item Có thể học decision boundary phức tạp
\end{itemize}

\textbf{Nhược điểm:}
\begin{itemize}
    \item Yêu cầu dataset lớn có nhãn (tốn kém!)
    \item Không phát hiện được anomaly mới (chưa thấy trong training)
\end{itemize}

\textbf{Trong nghiên cứu này:} LightAD (Decision Tree, SLFN)

\subsubsection{Semi-Supervised Learning (Học bán giám sát)}

Huấn luyện chủ yếu trên dữ liệu normal, có ít hoặc không có anomaly.

\textbf{Approaches:}
\begin{enumerate}
    \item \textit{One-class classification:} Học boundary xung quanh normal data (DeepSVDD)
    \item \textit{Few-shot learning:} Sử dụng vài mẫu anomaly để tinh chỉnh
\end{enumerate}

\textbf{Ưu điểm:}
\begin{itemize}
    \item Thực tế: Dữ liệu normal dồi dào, anomaly hiếm
    \item Có thể phát hiện novel anomaly (chưa thấy trong training)
\end{itemize}

\textbf{Nhược điểm:}
\begin{itemize}
    \item False positive rate cao hơn supervised
    \item Giả định training data "sạch" (không có hidden anomaly)
\end{itemize}

\textbf{Trong nghiên cứu này:} LogFormer (pre-train unsupervised + tune supervised), LogRAG (DeepSVDD + RAG)

\subsubsection{Unsupervised Learning (Học không giám sát)}

Không có nhãn; phát hiện anomaly dựa trên sự chênh lệch so với phân phối dữ liệu.

\textbf{Approaches:}
\begin{itemize}
    \item Clustering: Anomaly xa cluster center
    \item Density estimation: Anomaly có xác suất thấp
    \item Reconstruction: Anomaly có reconstruction error cao (Autoencoder)
\end{itemize}

\textbf{Ưu điểm:}
\begin{itemize}
    \item Không tốn chi phí gán nhãn
    \item Có thể khám phá anomaly type chưa biết
\end{itemize}

\textbf{Nhược điểm:}
\begin{itemize}
    \item False positive rate cao nhất
    \item Khó đánh giá nếu không có labeled test set
\end{itemize}

\textbf{Trong nghiên cứu này:} LightAD-KNN (có thể chạy unsupervised)

\begin{table}[ht]
\centering
\caption{So sánh các hướng tiếp cận theo mức độ giám sát}
\label{tab:supervision-comparison}
\begin{tabular}{lllll}
\hline
\textbf{Approach} & \textbf{Training Data} & \textbf{Ưu điểm} & \textbf{Nhược điểm} & \textbf{Use Case} \\
\hline
Supervised & Labeled (N+A) & High accuracy & Expensive labels & Known types \\
Semi-sup & Mostly normal & Detects novel & Higher FP & Rare anomalies \\
Unsupervised & Unlabeled & No label cost & Highest FP & Exploratory \\
\hline
\end{tabular}
\end{table}

\subsection{Các họ phương pháp}
\label{subsec:method-families}

\subsubsection{Classical Machine Learning}

\paragraph{Distance-Based Methods}

\textbf{K-Nearest Neighbors (KNN):} Phân loại anomaly dựa trên khoảng cách đến k láng giềng gần nhất.

\textbf{Thuật toán:}
\begin{enumerate}
    \item Cho một test sample $x$, tìm $k$ nearest neighbors trong training set
    \item Tính average distance $d_{avg}$
    \item Nếu $d_{avg} > \theta$ (threshold), classify as anomaly
\end{enumerate}

\textbf{Complexity:} $O(n^2)$ cho inference (brute-force), có thể giảm xuống $O(n \log n)$ với tree-based index.

\textbf{Ưu điểm:}
\begin{itemize}
    \item Đơn giản, interpretable (show nearest neighbors)
    \item Không cần training (instance-based learning)
    \item Hiệu quả với point anomaly
\end{itemize}

\textbf{Nhược điểm:}
\begin{itemize}
    \item Chi phí inference cao (cần tính distance với tất cả training samples)
    \item Nhạy cảm với feature scaling
    \item Curse of dimensionality (hiệu suất giảm trong không gian cao chiều)
\end{itemize}

\textbf{Trong nghiên cứu này:} LightAD-KNN

\paragraph{Tree-Based Methods}

\textbf{Decision Tree (DT):} Học các if-then rules để phân loại.

\textbf{Example rule:}
\begin{verbatim}
IF event_frequency("ERROR") > 10 AND 
   event_frequency("Timeout") > 5 
THEN anomaly = True
\end{verbatim}

\textbf{Complexity:} $O(n \log n)$ training, $O(\log n)$ inference (tree depth).

\textbf{Ưu điểm:}
\begin{itemize}
    \item \textbf{Highly interpretable} - có thể visualize cây quyết định
    \item Xử lý được mixed data types (categorical + numerical)
    \item Không cần feature scaling
    \item Inference rất nhanh
\end{itemize}

\textbf{Nhược điểm:}
\begin{itemize}
    \item Dễ overfit (single tree)
    \item Có thể không capture được complex patterns
\end{itemize}

\textbf{Trong nghiên cứu này:} LightAD-DT

\paragraph{Shallow Neural Networks}

\textbf{Single-Layer Feedforward Network (SLFN):} Neural network với một hidden layer duy nhất.

\textbf{Architecture:}
\[
\text{Input} \xrightarrow{W_1} \text{Hidden (ReLU)} \xrightarrow{W_2} \text{Output (Sigmoid)}
\]

\textbf{Complexity:} $O(n \cdot d \cdot h)$ với $d$ = input dim, $h$ = hidden units.

\textbf{Ưu điểm:}
\begin{itemize}
    \item Training nhanh hơn deep learning
    \item Có thể học non-linear boundaries
    \item Chi phí thấp
\end{itemize}

\textbf{Nhược điểm:}
\begin{itemize}
    \item Khả năng biểu diễn hạn chế (shallow)
    \item Ít interpretable hơn decision tree
    \item Có thể underfit với complex patterns
\end{itemize}

\textbf{Trong nghiên cứu này:} LightAD-SLFN

\subsubsection{Deep Learning}

\paragraph{Recurrent Networks (LSTM)}

\textbf{Core Idea:} Mô hình hóa temporal dependencies trong chuỗi log event.

\textbf{Approach:}
\begin{enumerate}
    \item Parse logs thành event sequences: $[e_1, e_2, ..., e_n]$
    \item Train LSTM để predict next event: $P(e_{t+1} | e_1, ..., e_t)$
    \item Anomaly: Next event không được dự đoán (probability thấp)
\end{enumerate}

\textbf{Key Papers:}
\begin{itemize}
    \item \textbf{DeepLog} \cite{du2017deeplog}: Đầu tiên áp dụng LSTM cho LogAD
    \item \textbf{LogAnomaly} \cite{meng2019loganomaly}: Kết hợp semantic + quantitative features
    \item \textbf{LogRobust} \cite{zhang2019logrobust}: Robust với unstable log data
\end{itemize}

\textbf{Ưu điểm:}
\begin{itemize}
    \item Capture temporal dependencies
    \item Hiệu quả với collective anomaly
    \item Xử lý variable-length sequences
\end{itemize}

\textbf{Nhược điểm:}
\begin{itemize}
    \item Training chậm (sequential)
    \item Vanishing gradient với long sequences
    \item Yêu cầu large training data
\end{itemize}

\paragraph{Transformer-Based Methods}

\textbf{Core Idea:} Sử dụng self-attention để capture long-range dependencies.

\textbf{LogFormer} \cite{guo2024logformer} (AAAI 2024) - \textit{Trong nghiên cứu này}:

\textbf{Architecture:}
\begin{enumerate}
    \item \textbf{Phase 1: Pre-training}
    \begin{itemize}
        \item Input: Unlabeled log sequences (100M logs từ BigLog corpus)
        \item Task: Masked Language Modeling (MLM)
        \item Output: Pre-trained Transformer
    \end{itemize}
    
    \item \textbf{Phase 2: Adapter Tuning}
    \begin{itemize}
        \item Input: Labeled log sequences (HDFS, BGL)
        \item Task: Binary classification (normal/anomaly)
        \item Method: Chỉ update adapter layers (1-5\% params)
        \item Output: Fine-tuned model for LogAD
    \end{itemize}
\end{enumerate}

\textbf{Ưu điểm:}
\begin{itemize}
    \item \textbf{Parallel training} - nhanh hơn LSTM
    \item \textbf{Long-range dependencies} - attention mechanism
    \item \textbf{Pre-training} - tận dụng unlabeled data
    \item \textbf{Parameter-efficient} - adapter tuning chỉ update 1-5\% tham số
\end{itemize}

\textbf{Nhược điểm:}
\begin{itemize}
    \item High GPU memory (self-attention là $O(n^2)$)
    \item Yêu cầu large corpus cho pre-training
    \item Ít interpretable (attention weights cần ML expertise)
\end{itemize}

\subsubsection{Large Language Models}

\paragraph{RAG-Based Methods}

\textbf{LogRAG} \cite{jiang2024lograg} (ISSRE 2024) - \textit{Trong nghiên cứu này}:

\textbf{Core Idea:} Kết hợp anomaly detection (DeepSVDD) với LLM reasoning (GPT-3.5/4) thông qua Retrieval-Augmented Generation.

\textbf{Architecture:}
\begin{enumerate}
    \item \textbf{Phase 1: DeepSVDD Anomaly Detection}
    \begin{itemize}
        \item Train one-class neural network trên normal logs
        \item Học hypersphere boundary xung quanh normal data
        \item Output: Anomaly score cho mỗi log session
    \end{itemize}
    
    \item \textbf{Phase 2: Retrieval-Augmented Generation}
    \begin{itemize}
        \item Với mỗi candidate anomaly:
        \begin{enumerate}
            \item Retrieve top-k similar log sessions (using FAISS index)
            \item Construct prompt: "Given examples: [Retrieved], is [Test] anomalous?"
            \item LLM reasoning: GPT-3.5/4 generates explanation
            \item Final decision: Combine DeepSVDD score + LLM reasoning
        \end{enumerate}
    \end{itemize}
\end{enumerate}

\textbf{Ưu điểm:}
\begin{itemize}
    \item \textbf{Explainable:} LLM cung cấp natural language explanation
    \item \textbf{Few-shot learning:} Retrieve similar examples
    \item \textbf{Semi-supervised:} DeepSVDD phase sử dụng unlabeled data
    \item \textbf{Domain knowledge:} LLM có thể reasoning về log semantics
\end{itemize}

\textbf{Nhược điểm:}
\begin{itemize}
    \item \textbf{High cost:} API calls cho mỗi test sample (GPT-3.5: \$0.002/1K tokens)
    \item \textbf{Latency:} Retrieval + LLM inference chậm (1-10 seconds/sample)
    \item \textbf{Hallucination risk:} LLM có thể generate plausible nhưng incorrect explanations
    \item \textbf{Dependency:} Phụ thuộc vào external API (OpenAI)
\end{itemize}

\subsection{So sánh ba phương pháp trong nghiên cứu}
\label{subsec:our-methods-comparison}

Bảng \ref{tab:methods-comparison} so sánh chi tiết ba phương pháp được đánh giá trong luận văn này.

\begin{table}[ht]
\centering
\caption{So sánh LightAD, LogFormer, và LogRAG}
\label{tab:methods-comparison}
\begin{tabular}{llll}
\hline
\textbf{Dimension} & \textbf{LightAD} & \textbf{LogFormer} & \textbf{LogRAG} \\
\hline
ML Family & Classical & Deep Learning & LLM \\
Architecture & KNN/DT/SLFN & Transformer & DeepSVDD + RAG \\
Supervision & Sup/Unsup & Semi-sup & Semi-sup \\
Training Data & 100K-1M & 10M-100M & 100K-1M \\
Training Time & Minutes & Hours-Days & Hours \\
Inference Time & <1ms (DT) & ~100ms & 1-10s \\
Interpretability & ⭐⭐⭐⭐ (DT) & ⭐⭐ (Attention) & ⭐⭐⭐⭐⭐ (NL) \\
Accuracy & Good & State-of-art & State-of-art \\
Cost & Very Low & Medium & High \\
\hline
\end{tabular}
\end{table}

\textbf{Key Insights:}
\begin{itemize}
    \item \textbf{LightAD:} Hiệu quả, nhanh, interpretable - phù hợp cho production với latency requirement
    \item \textbf{LogFormer:} Accuracy cao, tận dụng unlabeled data - phù hợp khi có large log corpus
    \item \textbf{LogRAG:} Explainability tốt nhất, few-shot - phù hợp khi cần natural language explanation cho SREs
\end{itemize}

\subsection{Metrics đánh giá}
\label{subsec:evaluation-metrics}

Do class imbalance cao (95-99\% normal logs), các metrics truyền thống như Accuracy không phù hợp. Chúng tôi sử dụng các metrics sau:

\subsubsection{Precision, Recall, F1-Score}

\begin{align}
\text{Precision} &= \frac{TP}{TP + FP} \label{eq:precision} \\
\text{Recall} &= \frac{TP}{TP + FN} \label{eq:recall} \\
\text{F1-Score} &= \frac{2 \cdot \text{Precision} \cdot \text{Recall}}{\text{Precision} + \text{Recall}} \label{eq:f1}
\end{align}

Trong đó:
\begin{itemize}
    \item $TP$ (True Positive): Anomaly được dự đoán đúng
    \item $FP$ (False Positive): Normal được dự đoán nhầm là anomaly
    \item $FN$ (False Negative): Anomaly bị bỏ sót (dự đoán là normal)
    \item $TN$ (True Negative): Normal được dự đoán đúng
\end{itemize}

\textbf{F1-Score} là harmonic mean của Precision và Recall, cân bằng giữa hai metrics này.

\subsubsection{PR-AUC (Area Under Precision-Recall Curve)}

Với imbalanced data, PR-AUC phù hợp hơn ROC-AUC \cite{saito2015precision}.

\textbf{Cách tính:}
\begin{enumerate}
    \item Vary threshold $\theta$ từ 0 đến 1
    \item Với mỗi $\theta$, tính Precision và Recall
    \item Plot PR curve: Recall (x-axis) vs Precision (y-axis)
    \item PR-AUC = Area under curve
\end{enumerate}

\textbf{Interpretation:} PR-AUC = 1.0 (perfect), Random classifier ≈ anomaly rate (e.g., 0.03 for 3\% anomaly).

\subsubsection{False Positive Rate (FPR)}

\begin{equation}
\text{FPR} = \frac{FP}{FP + TN}
\label{eq:fpr}
\end{equation}

\textbf{Importance:} Trong production, high FPR gây "alert fatigue" - SREs bỏ qua alerts do quá nhiều false alarms.

\textbf{Target:} FPR < 1\% thường được chấp nhận trong industry.

\subsection{Kết luận}
\label{subsec:anomaly-conclusion}

Phần này đã trình bày:
\begin{enumerate}
    \item Định nghĩa và phân loại bất thường (Point, Collective, Contextual)
    \item Các hướng tiếp cận theo supervision (Supervised, Semi-supervised, Unsupervised)
    \item Các họ phương pháp từ Classical ML đến Deep Learning và LLM
    \item Chi tiết về 3 phương pháp trong nghiên cứu (LightAD, LogFormer, LogRAG)
    \item Metrics đánh giá phù hợp với class imbalance
\end{enumerate}

Chương tiếp theo sẽ trình bày phương pháp nghiên cứu chi tiết, bao gồm experimental protocol, research questions, và pipeline thực nghiệm.
